\chapter{Roots and Weights}

\section{Motivating example: \(\mathfrak{su}(2)\)}

\subsection{Projective representations and simply-connected groups}

representations as linear maps on vectors

when we say the representation of a group is in GL or the one of an algebra is in End, we are saying that the group or algebra is represented as linear maps on vectors.

when we measure we normalize with resoect to the inner product, we need the norm of the vectors and thus the inner product...

all of the vectors on the same ray (same direction) are to be considered in the same state. it is a one dimensional subspace of the vector space. if this is what carachterizes the system, we need to really care only about the projected representation on the projective space, which is the space of rays. in qm symmetries should be realized as projective representattions of the group, rather than as linear representations. but we can always lift a projective representation to a linear one, by going to the universal cover of the group. for example, the projective representations of SO(3) are the linear representations of SU(2). this is why we can use SU(2) to describe spin, even if spin is not a representation of SO(3).

the mathematical key is the one saying that the a finite dimensional projective representation of a lie group on a finite dimensional complex vector space lifts to a unitary linear representation the covering group \(\overline{G}\), which is the unique simply connected Lie group with the same lie algebra as \(G\). hard to prove, absolutely non trivial.

projective representations \(\equiv\) unitary linear representations of the covering group \(\equiv\) anti hermitian linear representations of the lie algebra \(\equiv\) linear representations of the complexified lie algebra

In fact, for the special case of \(G=SO(3)\) and \(\mathfrak{g} = \mathfrak{su}(2)\), this math-result, combined with the “axiomatic” definition of quantum mechanics using the (projective) Hilbert space, can be viewed as predicting half-integer spin particles (i.e., linear representations of \(\mathfrak{su}(2)\) which are not \(SO(3)\) reps) in the quantum world: they arise from projective representations of the rotation symmetry \(SO(3)\), which we do not have in classical mechanics.

    [details on axiomatics of qm]

\subsection{Ladder Operators and Spin Representation}

\paragraph{Complexification and Hermitian basis.} from properties of complexification we can take complex linear combinations of the generators of the algebra, and we can choose them in a way that they have nice properties. since the original generators are anti-hermitian, we can take linear combinations that are hermitian. so we can diagonalize them. thinking of the adjoint rep in that basis, which is an endomorphism of the algebra, we can see that the adj rep of the generators is also an endomorphism of the algebra, and thus we can diagonalize it as well, since it will remain hermitian. this property will extend to all the other complex finite dimensional representations, and thus they will have a basis in which the generators are hermitian and thus diagonalizable.

\paragraph{Eigenvalues of \(d(J_3)\) from ladder operators.} we can pick an on basis thanks to the usual unitarity trick which allows us to make the representation unitary, and thus the generators hermitian. in the following we prove that for irreducible representations the dimension of the subspace of eigenvectors of \(d(J_3)\) with a given eigenvalue is one, and that the eigenvalues are equally spaced.
the ladder operators are not hermitian. after definingthem through the commutation relations, we can study how the eigenvalues of \(d(J_3)\) change when we apply the ladder operators (it moves up and down in the different eigenvector subspaces, the state changes but it remains an eigenvector). we can see that they change by one unit, and thus the eigenvalues are equally spaced (and thus the entire lie algebra moves among the invariant eigenvectors subspaces). this dont work in irreducible representations, since there do not exist trivial invariant subspaces.
if the representation is finite, there must exist a maximum range of eigenvalues, and thus there must be a highest weight vector, which is an eigenvector of \(d(J_3)\) that is annihilated by the raising operator. same for the minimum value (annihilated by the lowering operator).thus the representation space is made of a finite direct sum of eigenspaces of \(d(J_3)\) with equally spaced eigenvalues, and the ladder operators move us among them.

\paragraph{The Casimir operators and range of \(J_3\)-eigenvalues.} to further restrict the ...
the casimir operator commutes with all the generators, thus it is a multiple of the identity in an irreducible representation. we can apply schurs lemma to infer that the eigenvalue of the casimir operator is a constant in an irreducible representation.
if we apply the obtained formulae to the basis of eigenvectors of \(d(J_3)\), we can find a relation between the eigenvalue of the casimir operator and the range of the eigenvalues of \(d(J_3)\). this allows us to find the possible values of the eigenvalues of \(d(J_3)\) and thus to classify all the irreducible representations. the proprierties we derive constrain the spacing of the eigenvalues to be integer or half-integer (and \(>0\)), and thus we can label the irreducible representations by a non-negative integer or half-integer \(j\), which is the maximum eigenvalue of \(d(J_3)\) (the highest weight). the dimension of the representation is then \(2j+1\).

\paragraph{Dimensionality of eigenspaces.} we want to argue that all the eigenspaces of \(d(J_3)\) are one-dimensional. the span of the maximum eigenvectors and consecutive applications of the lowering operator is an invariant subspace, thus it must be the entire representation space, since the representation is irreducible. thus all the eigenvectors of \(d(J_3)\) are obtained by applying the lowering operator to the maximum eigenvector, and thus they are all one-dimensional (we never get to the same eigenvalue twice, since the eigenvalues are equally spaced and we have a finite number of them).

\paragraph{Normaliation and numerical coefficients.} [...] we find how to completely specify the representation by giving the value of \(j\) and the normalization of the ladder operators. we can choose a normalization such that the action of the ladder operators on the eigenvectors is given by a simple formula, which is very convenient for computations.

\section{Reformulation with Roots and Weights}

 [...]

\section{Killing Form and Cartan Subalgebra}

For a general Lie algebra \(\mathfrak{g}\), we first need the analogue of the subalgebra spanned by \(J_3\) in \(\mathfrak{sl}(2,\mathbb{C})\)  which we choose to diagonalize.

\begin{definition}[Killing form]
    [...]
\end{definition}

It takes two elements of the algebra and gives a number, thus it is a bilinear form on the algebra. since we are taking a trace of endomorphisms, we now there exist a basis invariant way to define it, and thus it is invariant under the adjoint action of the algebra on itself. this is a very important property, since it allows us to use the killing form to define an inner product on the algebra, which is invariant under the adjoint action. this will allow us to define a notion of orthogonality and thus to define a notion of diagonalization for the adjoint representation, which is crucial for the classification of representations.

The symmetric property comes from the fact that the trace of a product of endomorphisms is invariant under cyclic permutations, thus we can swap the two elements without changing the value of the killing form. the invariance property comes from the fact that the killing form is defined in terms of the adjoint representation, which is invariant under the adjoint action.

Given a basis for the lie algebra, given by the generators, we can write the killing form in terms of the structure constants of the algebra, with structure constants defined by the commutation relations of the generators. we know tha adjoint representation is given by the structure constants, thus we can write the killing form in terms of the structure constants. this allows us to compute the killing form explicitly for any given lie algebra, once we know its structure constants. Thus it is written in the same basis.

\[
    K_{mn} = \dots
\]

in particular for \(\mathfrak{g} = \mathfrak{su}(2)\) in the previous antihermitian basis we have \(K_{mn} = -\delta_{mn}\) and for the hermitian basis for the complexification \(\mathfrak{sl}(2,\mathbb{C})\)  we have \(K_{mn} = \delta_{mn}\).

There is an ambiguity when it comes to subalgebras: there are a priori two pairings we can define on a subalgebra, the one given by the killing form of the original algebra and the one given by the killing form of the subalgebra (which is a Lie algebra iself). in general they are not the same, thus we need to choose which one to use. whats relevant for us is when they are the same: in any case in which the subalgebra is an invariant subspace of the adjoint then we have the same killing form, since the killing form is defined in terms of the adjoint representation. thus we will always choose a subalgebra that is an invariant subspace of the adjoint, so that we can use the same killing form to define an inner product on it.
thus for a semisimple lie algebra, we have a unique killing form, made by the direct sum of the killing forms of the simple components.
another relevant case is when the subalgebra is real (from the complexification) then the real basis can be used to define the complex basis, therefore structure constants and the pairing matrix in the basis are the same, thus the following corollary holds:

\begin{corollary}
    Let \(\mathfrak{h}\) be a real Lie algebra and \(\mathfrak{g} = \mathfrak{h}_{\mathbb{C}}\) its complexification. Then the killing form of \(\mathfrak{h}\) and the killing form of \(\mathfrak{g}\) coincide on \(\mathfrak{h}\):
    \[
        K^{\mathfrak{h}} = K^{\mathfrak{g}} \big|_\mathfrak{h}.
    \]
\end{corollary}

\begin{proposition}
    The Killing form on \(\mathfrak{g}\) satisfies \(\forall\) ...
    \[
        \dots
    \]
    symmetric pairing - invariant
    their sum is zero
\end{proposition}
\begin{proof}
    [...]
\end{proof}

The notion of “invariance” becomes intuitive when we inspect the interplay between the Killing form on \(\mathfrak{g}\) and the adjoint representation of the Lie group \(G\) with algebra \(\mathfrak{g}\). Namely, from Corollary 2.3.3 we have for all \(X,Y,Z \in \mathfrak{g}\):
\[
    \dots
\]
from which we can see that the infinitesimal transformation \(\mathrm{ad}_{\exp(Z)}\) does not change the Killing form of any two elements of the algebra, thus it is invariant under the adjoint action of the group on the algebra. this is a very important property, since it allows us to use the killing form to define an inner product on the algebra, which is invariant under the adjoint action.

\begin{example}[...]
    [...]
\end{example}

[...]

\begin{theorem}
    A Lie algebra is semi-simple if and only if the Killing form is a non-degenerate bilinear form, i.e., if \(K(X,Y)=K(Y,X)=0\)  for all \(Y \in \mathfrak{g}\), then \(X = 0\).
\end{theorem}

While this theorem can be clearly stated, and is important for us in the following, its proof utilizes concepts of ideals and radicals of Lie algebras, which we have not introduced. This is a bit hard to prove so we wont.

\begin{theorem}
    Let \(G\) be a compact real Lie group and \(\mathfrak{g}\) its Lie algebra. Then its Lie algebra has to be \textbf{reductive}:
    \[
        \mathfrak{g} \equiv \mathfrak{g}^{(ab)} \oplus \mathfrak{g}^{(ss)},
    \]
    where \(\mathfrak{g}^{(ab)}\) is an abelian Lie algebra and \(\mathfrak{g}^{(ss)}\) is a semisimple Lie algebra. Furthermore, the Killing form is negative-semi-definite: ...

    Conversely, if \(\mathfrak{g}\) is a real Lie algebra with negative-definite Killing form (then \(\mathfrak{g}\) must be semi-simple), then any connected Lie group with algebra \(\mathfrak{g}\) is compact. If a real Lie algebra \(\mathfrak{g}\) has positive-definite Killing form, then \(\mathfrak{g} = {0}\).
\end{theorem}

Thus the Killing form relates to the semisimplicity of the algebra, and thus to the compactness of the group. In particular, for a compact real Lie group, the killing form is negative-semi-definite, and thus it is negative-definite on the semisimple part of the algebra, and zero on the abelian part.

The relevance of these results to us is highlighted in the following way. Consider first a real semi-simple algebra \(\mathfrak{g}\) with compact \(G_0\). Pick an \(\mathbb{R}\)-basis \(\{T_i\}\) such that
\[
    K_{ij} = K(T_i, T_j) = \text{diag}\left(-1, \dots, -1\right).
\]
which is always possible by Gram–Schmidt. One particularly nice property of this basis is that the structure constants
\[
    [T_i, T_j] = \mathrm{ad}_{T_i}(T_j) = c_{ij}^k T_k,
\]
are constrained by the invariance of the Killing form
\[
    \dots
\]
thus using the invariance and the symmetry of the Killing form we can derive a constraint on the structure constants, which is a very useful property for computations: \uline{they are \textbf{completely} \textbf{antisymmetric} in the three indices}, thus they are the structure constants of a Lie algebra with respect to an orthonormal basis. We already know that the structure constants were antisymmetric in the first two indices (due to the definition through commutators), but now we have a stronger constraint, which is that they are completely antisymmetric in all three indices.

Passing to the complexification, we can define a \(\mathbb{C}\)-basis of \(\mathfrak{g}_\mathbb{C}\), \(\{J_m = i T_m\}\), which by Corollary 3.3.2 is orthonormal, i.e., \(K^{\mathfrak{g}_\mathbb{C}}(J_m, J_n) = K^{\mathfrak{g}_\mathbb{C}}(i T_m, i T_n) = -K^{\mathfrak{g}_\mathbb{C}}(T_m, T_n) = -K(T_m, T_n) = \delta_{mn}\). In this new basis, the adjoint rep \(\mathrm{ad}\) of \(\mathfrak{g}_\mathbb{C}\) satisfies
\[
    \dots
\]
which is an Hermitian matrix. We now know that in this basis we can diagonalize the adjoint representation, since we have found a basis in which the adjoint representation is Hermitian. So \(\mathrm{ad}_{J_l} \in \mathrm{End}(\mathfrak{g}_\mathbb{C}) \) for the particular generators \(J_l\) of the complexified algebra can be diagonalized.

Given the many details / proofs we have left out in arriving at this statement, we summarize the main result in the following:
\begin{corollary}
    Let \(\mathfrak{g}\) be a complex semi-simple Lie algebra. Then there exists elements \(X \in \mathfrak{g}\) such that \(\mathrm{ad}_X \in \mathrm{End}(\mathfrak{g})\) is diagonalizable.
\end{corollary}
This is also a sort of justification for the habit of physicists to use infinitesimal transformations with complex coefficients... you can find the real.... ?

In the following, we will mostly work on complex algebras, so \(\mathfrak{g}\) will denote a complex semi-simple Lie algebra unless otherwise stated.

\subsubsection{Cartan subalgebra}

\begin{definition}
    We have two statements
    \begin{enumerate}[label=(\roman*)]
        \item toroidal...
        \item cartan subalgebra...
    \end{enumerate}
\end{definition}

In the example \(\mathfrak{g} = \mathfrak{su}(2)_\mathbb{C} = \mathfrak{sl}(2; \mathbb{C})\), we alredy saw that \(J_3\) spans a Cartan subalgebra (since the generators do not commute we can take one of them). However it is also clear that any of the other generators \(J_{1/2}\) would equally span a Cartan algebra. So the \uline{Cartan is no unique, but basically “a basis change away”}.

This is made precise in the following
\begin{theorem}
    Let \(\mathfrak{h}_1, \mathfrak{h}_2 \subset \mathfrak{g}\) be two Cartan subalgebras of a complex semi-simple Lie algebra \(\mathfrak{g}\). Then there exists an element in any group \(G\) with such algebra that maps one Cartan subalgebra to the other by the adjoint action:
    \[
        \mathfrak{h}_1 = \mathrm{Ad}_\mathfrak{g} (\mathfrak{h}_2).
    \]
    Consequently, \(\mathrm{dim}(\mathfrak{h}_1) = \mathrm{dim}(\mathfrak{h}_2)\).
\end{theorem}
This is in a sense just a rotation, since the Killing form is left invariant by the adjoint action, thus it is a rotation in the space of generators. Thus all Cartan subalgebras are equivalent.

\begin{definition}
    The \textbf{rank} of a complex semi-simple Lie algebra \(\mathfrak{g}\) is the dimension of its Cartan subalgebra(s) \(\mathfrak{h}\):
    \[
        \mathrm{rank}(\mathfrak{g}) := \mathrm{dim}(\mathfrak{h}).
    \]
\end{definition}
Thus \(\mathfrak{su}(2)_\mathbb{C} = \mathfrak{sl}(2; \mathbb{C})\) has rank 1; more generally, \(\mathfrak{su}(n)_\mathbb{C} = \mathfrak{sl}(n; \mathbb{C})\) has rank \(n - 1\); \(\mathrm{rank}(\mathfrak{so}(2n)_\mathbb{C}) = \mathrm{rank}(\mathfrak{so}(2n + 1)_\mathbb{C}) = n\).

The name “toral/toroidal” algebra comes from the following. By definition, a toral algebra \(\mathfrak{h} \subset \mathfrak{g}\) is Abelian. If we pick a compact real form \(\mathfrak{k} \subset \mathfrak{g}\), then \(\mathfrak{h} \cap \mathfrak{k}\) is a real toral algebra, which exponentiates into a compact Abelian subgroup (these elements commute), i.e., of the form \(\mathrm{U}(1)^n \equiv (\mathbb{S}^1)^n \) which is topologically a torus. A Cartan subalgebra exponentiates into a “maximal torus”, which is a common name in the literature. Just as their algebras, maximal tori \(\mathcal{H}_{1,2}\) are conjugate to each other:
\[
    \exists g \in G \,:\, g\mathcal{H}_1g^{-1} = \mathcal{H}_2.
\]

\section{Roots and Root Decomposition}

In the following, we assume that we have fixed a Cartan algebra \(\mathfrak{h} \subset \mathfrak{g}\); other choices would yield identical results given that they differ by a basis change \(\mathrm{Ad}_g \in G \to \mathrm{GL}(\mathfrak{g})\). Since by definition all elements of \(\mathfrak{h}\) commute, i.e.
\[
    \forall H, H' \in \mathfrak{h}, \quad [\mathrm{ad}_H, \mathrm{ad}_{H'}] = \mathrm{ad}_{[H, H']} = 0,
\]
we can diagonalize \(\mathrm{ad}_H\) for all \(H \in \mathfrak{h}\) simultaneously, i.e., we can find a basis of \(\mathfrak{g}\) which consists of (simultaneous) eigenvectors \(X\) of all \(\mathrm{ad}_H\), whose eigenvalue we denote by \(\lambda(H) \in \mathbb{C}\):
\[
    \implies \mathrm{ad}_H(X) = [H, X] = \lambda(H) X.
\]

Such eigenvalues satisfy for any \(H_1, H_2 \in \mathfrak{h}\):
\[
    \begin{aligned}
        \mathrm{ad}_{H_1 + H_2}(X) & = \mathrm{ad}_{H_1}(X) + \mathrm{ad}_{H_2}(X) = \lambda(H_1) X + \lambda(H_2) X = (\lambda(H_1) + \lambda(H_2)) X, \\
        \mathrm{ad}_{c H_1}(X)     & = c \mathrm{ad}_{H_1}(X) = c \lambda(H_1) X, \quad c \in \mathbb{C},
    \end{aligned}
\]
So every (simultaneous) eigenvector \(X \in \mathfrak{g}\) defines a linear functional \(\lambda \,:\, \mathfrak{h} \to \mathbb{C}\), \(H \mapsto \lambda(H)\), i.e., \(\lambda \in \mathfrak{h}^\vee\) is an element of the vector space dual to the Cartan algebra.

For any \(\alpha \in \mathfrak{h}^\vee\), we can then define the eigenspace
\[
    \mathfrak{g}_{\alpha} := \{ X \in \mathfrak{g} \,:\, \mathrm{ad}_H(X) = [H, X] = \alpha(H) X, \forall H \in \mathfrak{h} \},
\]
i.e., the set of eigenvectors whose eigenvalues agree with the linear functional \(\alpha\). Note that, as defined here, \(\mathfrak{g}_{\alpha}\) is a vector space, but a priori not a Lie algebra. Clearly, a generic \(\alpha \in \mathfrak{h}^\vee\) will not be the eigenvalues of \(\mathrm{ad}_H\) for any \(H \in \mathfrak{h}\), so \(\mathfrak{g}_{\alpha} = \{0\}\) generically. Furthermore, since \(\mathfrak{h}\) is the maximal Abelian subalgebra, it is easy to see that \(\mathfrak{h} = \mathfrak{g}_0\). \uline{So this is a vector space but not a Lie algebra in general}.

\begin{definition}[Root system]
    [...]
\end{definition}

The set of all dual vector as appears ...

\begin{lemma}
    \(\mathcal{R}\) is a finite set.
\end{lemma}
\begin{proof}
    This is obvious from the fact that eigenspaces \(\mathfrak{g}_{\alpha}\) associated to different eigenvalues of a diagonalizable endomorphism are linearly independent. Since \(\mathfrak{g}\supset \mathfrak{g}_{\alpha}\) is finite dimensional, there can only be finitely many \(\alpha\) such that \(\mathfrak{g}_\alpha \neq \{0\}\).
\end{proof}

Let us revisit \(\mathfrak{g}=\mathfrak{sl}(2; \mathbb{C})\). We fix \(\mathfrak{h} = \{ \mu J_3 \,|\, \mu \in \mathbb{C} \} \subset \mathfrak{sl}(2; \mathbb{C})\) as the Cartan algebra. Given the commutation relation \([J_3,\,J_{\pm}] = \pm J_{\pm}\), we see that \(\mathrm{ad}_{J_3}\) only has \(\pm 1\) as non-zero eigenvalues. So \(R = \{ \alpha_+ \equiv \alpha, \alpha_- = -\alpha \} \subset \mathfrak{h}^\vee\) with \(\alpha : \mathfrak{h} \to \mathbb{C}\), \(J_3 \mapsto 1\), and \(\mathfrak{g}_{\alpha_\pm} = \{ \mu J_\pm \,|\, \mu \in \mathbb{C} \}\). Notice how, as a vector space, \(\mathfrak{sl}(2; \mathbb{C}) = \mathfrak{h} \oplus \mathfrak{g}_{\alpha_+} \oplus \mathfrak{g}_{\alpha_-} \). These features, together with a few other ones that are easy to check explicitly for \(\mathfrak{sl}(2; \mathbb{C})\), are the content of the following theorem.

\begin{theorem}[Root decomposition]
    We fix \(\mathfrak{h} \subset \mathfrak{g}\) as the Cartan subalgebra, and \(R\) the root system of \(\mathfrak{g}\) with \(\mathfrak{h} = \mathfrak{g}_0\).
    \begin{enumerate}[label=(\roman*)]
        \item We can decompose the Lie algebra as a vector space as
              \[
                  \mathfrak{g} = \mathfrak{h} \oplus \bigoplus_{\alpha \in \mathcal{R} \Subset} \mathfrak{g}_\alpha,
              \]
              as the cartan subalgebra plus the root spaces. This is called the root decomposition of the algebra. It is not a sum of algebras, since the root spaces are not closed under the Lie bracket, they are just vector spaces.
        \item If we take two such elements ... then the commutator lives in the space associated to the sum of the two roots ... and the brackets must be closed...
        \item If we have two roots that do not sum to zero ...
        \item For any root \(\alpha\), one can take the killing form restricted to the root space \(\mathfrak{g}_\alpha\) and the root space \(\mathfrak{g}_{-\alpha}\), and this gives a non-degenerate pairing between these two spaces, thus they have the same dimension.
    \end{enumerate}
\end{theorem}
\begin{proof}
    \begin{enumerate}[label=(\roman*)]
        \item Follows directly from \(\{\mathrm{ad}_H \,|\, H \in \mathfrak{h}\}\) being simultaneously diagonalizable, so \(\mathfrak{g}\) decomposes into the direct sum of the eigenspaces, which are by definition the root subspaces \(\mathfrak{g}_{\alpha}\).
        \item Direct computation gives for any \(H \in \mathfrak{h}\):
              \[
                  \dots
              \]
        \item We can rewrite (ii) also as \(\mathrm{ad}_{E_{\alpha}} (E_{\beta}) \in \mathfrak{g}_{\alpha + \beta}\). So, for \(E_{\gamma} \in \mathfrak{g}_{\gamma}\) we have
              \[
                  \mathrm{ad}_{E_{\alpha}} \circ \mathrm{ad}_{E_{\beta}} = \sum_{\alpha, i} \lambda_{\alpha}^{(i)} E_\alpha^{(i)} \in \mathfrak{g}_{\alpha + \beta + \gamma},
              \]
              So if \(\alpha + \beta \neq 0\) then \(E_{\gamma} \notin \mathfrak{g}_{\alpha + \beta + \gamma}\), which implies that \(\lambda_{\gamma}^{(i)} = 0\): all of the diagonal elements of the matrix representing this endomorphism in the root-eigenbasis are zero.

              Then, unless \(\alpha + \beta = 0\), \(\mathrm{ad}_{E_{\alpha}} \circ \mathrm{ad}_{E_{\beta}} \) maps any basis vector \(T\) which by (i) can be chosen to be an eigenvector with certain eigenvalue \(\gamma \in \mathcal{R}\), to a linear combination of basis vectors that are eigenvectors with eigenvalues \(\alpha + \beta + \gamma \neq \gamma\). This again implies that the diagonal entries for the matrix of \(\mathrm{ad}_{E_{\alpha}} \circ \mathrm{ad}_{E_{\beta}} \) in the root-eigenbasis must be all 0, so \(\Tr(\mathrm{ad}_{E_{\alpha}} \circ \mathrm{ad}_{E_{\beta}}) = K(E_{\alpha}, E_{\beta}) = 0\).
        \item By (iii), the only elements \(0 \neq Y \in \mathfrak{g}\) such that \(K(X, Y) \neq 0\) for an element \(0 \neq X \in \mathfrak{g}_\alpha \subset \mathfrak{g}\) must be in \(\mathfrak{g}_{-\alpha} \subset \mathfrak{g}\). Since we assumed the Lie algebra to be semi-simple, then its Killing form \(K\) is non-degenerate on \(\mathfrak{g}\): there must exist such a \(0 \neq Y \in \mathfrak{g}_{-\alpha}\). So \(K\) provides a non-degenerate pairing \(\mathfrak{g}_\alpha \times \mathfrak{g}_{-\alpha} \to \mathbb{C}\). Thus, if \(\mathfrak{g}_\alpha\) is non-trivial, also \(\mathfrak{g}_{-\alpha} \neq \{0\}\), hence \(-\alpha \in \mathcal{R}\).
    \end{enumerate}
\end{proof}

Thus if we take commutators of vectors from different groups, we get vectors in the group associated to the sum of the roots. In particular, if we take commutators of vectors in the same group, we get vectors in the group associated to twice the root. So if \(2\alpha \notin \mathcal{R}\) then \(\mathfrak{g}_{2\alpha} = \{0\}\) and thus \([\mathfrak{g}_\alpha, \mathfrak{g}_\alpha] = \{0\}\).\TODO{check}

This is the single most important result about semi-simple Lie algebras. Everything we will learn further about their structure, including the classification and constructing their representations, builds on this decomposition.

\begin{definition}
    Using the fact that \(K|_{\mathfrak{h}}\) is a non-degenerate symmetric pairing, we have a vector space isomorphism \(\kappa \,:\, \mathfrak{h}^\vee \xrightarrow{\simeq} \mathfrak{h}\). This gives us a non trivial notion of duality, which is invertible:
    \[
        \begin{aligned}
            \kappa^{-1} \,:\, & \mathfrak{h} \xrightarrow{\simeq} \mathfrak{h}^\vee, \\
                              & H \mapsto \kappa(H, \cdot).
        \end{aligned}
    \]
    \begin{enumerate}
        \item Let \(\alpha \in \mathfrak{h}^\vee\), then we can denote \(H_{\alpha} := \kappa(\alpha)\). It satisfies \(K(H_{\alpha}, H_{\beta}) = \alpha(H_{\beta})\) for all \(\alpha, \beta \in \mathfrak{h}^\vee\).
        \item ...
    \end{enumerate}
\end{definition}

For a better intuition about it, let us fix a basis \(\{H_i\}\) of \(\mathfrak{h}\), in which \(K|_{\mathfrak{h}}\) has the matrix \(\kappa_{ij} = K(H_i, H_j)\). In the dual basis \(\{\theta_i\}\) of \(\mathfrak{h}^\vee\) (i.e., \(\theta_i(H_j) = \delta_{ij}\)), we have for all \(\alpha \in \mathfrak{h}^\vee\) the equation
\[
    H_\alpha = \kappa(\alpha_i \theta_i) = \kappa_{ij} \alpha_i H_j,
\]
where \(\kappa_{ij} \equiv (\kappa^{-1})_{ij}\) is the inverse matrix, which also encodes the pairing on \(\mathfrak{h}^\vee\):
\[
    \begin{aligned}
        (\alpha, \beta) = K(H_\alpha, H_\beta) = K(\kappa^{ij}\alpha_i H_j, \kappa^{lm}\beta_l H_m) = \alpha_i \beta_l \kappa^{ij} \kappa^{lm} K(H_j, H_m) \\
        = \alpha_i \beta_l \kappa^{ij} \kappa^{lm} \kappa_{jm} = \alpha_i \beta_l \kappa^{ij} \delta^l_{\ j} = k^{ij} \alpha_i \beta_j.
    \end{aligned}
\]

For \(\mathfrak{g} = \mathfrak{sl}(2; \mathbb{C})\), we have used for \(\mathfrak{h}\) the basis \(H = J_3\), with \(K(J_3, J_3) = 2\) and
\[
    [J_3, X] = \alpha(J_3) X, \quad [J_e, J_\pm] = \pm J_\pm,
\]
and the roots \(\alpha^{\pm} : J_3 \mapsto \pm 1\) associated with eigenvectors \(E_{\pm\alpha} = J_{\pm}\):
\[
    \mathfrak{g}_{\alpha^{\pm}} = \{ \mu J_{\pm} \,|\, \mu \in \mathbb{C} \} = \{ \mu E_{\pm\alpha} \,|\, \mu \in \mathbb{C} \}.
\]
Then:
\begin{enumerate}
    \item \(H_{\alpha} = \kappa(\alpha) = \lambda J_3 \in \mathfrak{h}\), is determined by the condition...

          If we choose a different basis for the Cartan subalgebra, for example \(J_1\) instead of \(J_3\), then we would have a different root system. However, the root system is basically the same...
    \item \(\alpha(J_3)=1\), they form a dual pair of basis for...
\end{enumerate}

We found the roots of the root system identifying the eigenvalues of the adjoint representation of the Cartan subalgebra...
We restrict the killing form to the cartan subspace and use it to define a non-degenerate pairing (given by a non degenerate matrix) between the root spaces associated to opposite roots, which allows us to find the dual vector \(H_{\alpha}\) associated to each root \(\alpha\) (it gives an isomorphism between \(\mathfrak{h}^\vee\) and \(\mathfrak{h}\)). We can use this matrix \(\kappa\) to identify vectors in the Cartan subalgebra with linear functionals on it. We can use this pairing as a geometric euclidean pairing on the root space, which allows us to define angles and lengths of roots (through scalar products). We can use the root decomposition to find the commutation relations among the generators of the algebra, which are constrained by the root structure.

Let us discuss some further properties of the root decomposition, which are important for the classification of representations.

\begin{lemma}
    Let \(E_+ \in \mathfrak{g}_{\alpha}\) and \(E_- \in \mathfrak{g}_{-\alpha}\), then
    \[
        [E_+ ,\, E_-] = K(E_+,E_-) = H_{\alpha}.
    \]
\end{lemma}
\begin{proof}
    [...]
\end{proof}

\begin{lemma}
    Let \(E_+ \in \mathfrak{g}_{\alpha}\) and \(E_- \in \mathfrak{g}_{-\alpha}\), then
    \[
        [E_+ ,\, E_-] = K(E_+,E_-) H_{\alpha},
    \]
    for \(\alpha \in \mathcal{R}\), where \(H_{\alpha} = \kappa(\alpha)\) is the dual vector associated to \(\alpha\) through the Killing form.
\end{lemma}
\begin{proof}
    \[
        H:=\dots
    \]
    \[
        K\big|_{\mathfrak{h}} = \dots
    \]
    This means that by leaving the second argument of the Killing form free, we can identify it with the dual vector \(H_{\alpha}\) associated to \(\alpha\) through the Killing form...
    \[
        H = \kappa(\kappa^{-1}(H)) =\dots
    \]
\end{proof}

\begin{lemma}
    Let \(\alpha \in \mathcal{R} \subset \mathfrak{h}^\vee\) be a root \(\mathfrak{g}\). Then \((\alpha, \alpha) = K(H_{\alpha}, H_{\alpha}) \neq 0\).
\end{lemma}
\begin{proof}
    Let us make a futher assumption: the basis of the cartan subalgebra can diagonalize ... such that \(\mathrm{ad}_{H_{\alpha}}\) is Hermitian (highly non trivial and we cannot prove it at this point, but it is true). Then we can compute
    \[
        K(H_{\alpha}, H_{\alpha}) = \Tr(\mathrm{ad}_{H_{\alpha}} \circ \mathrm{ad}_{H_{\alpha}}) = \sum_{\beta \in \mathcal{R}}  \beta(H_{\alpha})^2 \dim(\mathfrak{g}_\beta),
    \]
    we knew that we can diagonalize \(\mathrm{ad}_{H_{\alpha}}\), but knowing that it is Hermitian allows us to conclude that its eigenvalues are real, thus \(\beta(H_{\alpha})^2 \geq 0\) for all \(\beta \in \mathcal{R}\) (also the dimension of the subspaces are positive), so the sum is positive, thus \(K(H_{\alpha}, H_{\alpha}) > 0\).

    Thus assuming the hermiticity let us define the ``length'' of a root \(\alpha\) as \(|\alpha| := \sqrt{(\alpha, \alpha)} = \sqrt{K(H_{\alpha}, H_{\alpha})}\). We can now divide by \((\alpha, \alpha)\).
\end{proof}

\begin{lemma}
    Let \(J_+^{(\alpha)}\) ...
    \[
        K(J_+^{(\alpha)}, J_-^{(\alpha)}) = \frac{2}{(\alpha, \alpha)}.
    \]
    We needed the previous lemma to be able to divide by \((\alpha, \alpha)\) and thus to define the normalized generators \(J_{3}^{(\alpha)}\) associated to the root \(\alpha\):
    \[
        J_3^{(\alpha)}\coloneqq \frac{H_{\alpha}}{(\alpha, \alpha)}.
    \]
    Then...

    The claim is that we have found a subalgebra of \(\mathfrak{g}\) isomorphic to \(\mathfrak{sl}(2; \mathbb{C})_{\alpha}\) associated to each root \(\alpha\), which is generated by the three elements \(J_+^{(\alpha)}\), \(J_-^{(\alpha)}\) and \(J_3^{(\alpha)}\). This is called the \textbf{root subalgebra} associated to the root \(\alpha\).
\end{lemma}
\begin{proof}
    By direct computation
    \[
        \dots
    \]
    Evaluating \(\alpha\) on its dual vector \(H_{\alpha}\) we get \(\alpha(H_{\alpha}) = K(H_{\alpha}, H_{\alpha}) = (\alpha, \alpha)\) is just the norm; the commutator among \(H_{\alpha} \) and \(J_{\pm}^{(\alpha)}\) is just the eigenvalue of \(H_{\alpha}\) on the root space \(\mathfrak{g}_{\pm\alpha}\), since the latter ar eigenvalues of \(\mathrm{ad}_{H_{\alpha}}\) by definition.

    We have in the end found the commutation relations of \(\mathfrak{sl}(2; \mathbb{C})\) for the generators \(J_3^{(\alpha)}\), \(J_{\pm}^{(\alpha)}\).
\end{proof}

Every root of a semisimple lie algebra defines a subalgebra isomorphic to \(\mathfrak{sl}(2; \mathbb{C})\), which is generated by the three elements \(J_3^{(\alpha)}\), \(J_{\pm}^{(\alpha)}\). This is a very important result, since it allows us to describe the lie algebra in terms of the representations of this subalgebra; in particular, taking the adjoint representation of the lie algebra, we can restrict it to the root subalgebra (spanned by \(J_+^{\alpha}\), \(J_-^{\alpha}\) and \(J_3^{\alpha}\)), and thus we can decompose the algebra into irreducible representations of \(\mathfrak{sl}(2; \mathbb{C})\), on the space of endomorphisms of the algebra itself:
\[
    \mathrm{ad}^{(\mathfrak{g})} \big|_{\mathfrak{sl}(2; \mathbb{C})_{\alpha}} \,:\, \mathfrak{sl}(2; \mathbb{C})_{\alpha} \to \mathrm{End}(\mathfrak{g}).
\]

\begin{proposition}
    Let \(\alpha \in \mathcal{R}\) be a root of a complex semi-simple Lie algebra \(\mathfrak{g}\). Then let the root subalgebra \(\mathfrak{sl}(2; \mathbb{C})_{\alpha} = \mathrm{span}_{\mathbb{C}}\{J_3^{(\alpha)}, J_+^{(\alpha)}, J_-^{(\alpha)}\}\) associated to \(\alpha\). Then the subspace
    \[
        V\coloneqq\dots
    \]
    this is a vector space which defines an irreducible representation of \(\mathfrak{sl}(2; \mathbb{C})_{\alpha}\) under the adjoint action (induced by the restriction on the adjoint representation on the subspace): it is an \(\mathfrak{sl}(2; \mathbb{C})\) irreducible representation under
    \[
        \left\{\left. \mathrm{ad}^{(\mathfrak{g})}_{J_\pm^{(\alpha)}}\right|_V , \left. \mathrm{ad}^{(\mathfrak{g})}_{J_3^{(\alpha)}}\right|_V \right\}.
    \]
\end{proposition}

\begin{proof}
    First recall... spin... It is carachterized by
    \[
        d(J_3)\ket{j,m} = \dots
    \]
    Now, since \(\mathfrak{g}_{\alpha}\) are orthogonal to each other (except for \(\pm \alpha\)), there can only be a finite set of \(m \in \mathbb{Z} \backslash \{0\}\) for which \(\mathfrak{g}_{m \alpha}\) is non-trivial, since \(V \subset \mathfrak{g}\) must be finite dimensional.

    We want to show that \(V\) is invariant under \(\left\{\left. \mathrm{ad}^{(\mathfrak{g})}_{J_\pm^{(\alpha)}}\right|_V , \left. \mathrm{ad}^{(\mathfrak{g})}_{J_3^{(\alpha)}}\right|_V \right\}\), in order to prove that it is a representation space of \(\mathfrak{sl}(2; \mathbb{C})_{\alpha}\).
    \begin{itemize}
        \item \(\dots \)
        \item \(\dots \)the relevant cases are with \(m=\pm 1\) since in the definition of \(V\) we are excluding \(0\) and if we act with \(J_\pm^{(\alpha)}\) we get \(\pm 1\) in the eigenvalue, so we need to check the cases for \(m = \pm 1\) to be sure of the result...
    \end{itemize}
    This shows that this operators restricted on \(V\) are endomorphisms of \(V\), thus \(V\) is a representation space of \(\mathfrak{sl}(2; \mathbb{C})_{\alpha}\).

    From the first bullet point we know that the basis vectors of \(V\) are eigenvectors of \(\mathrm{ad}_{J_3^{(\alpha)}}\) with eigenvalues \(m \in \mathbb{Z} \backslash \{0\}\). Furthermore we know from this first point that the \(J_3^{(\alpha)}\)-eigenbasis in \(V\) has integer eigenvalues. If the subspace \(V\) splits in a direct sum of subspaces
    \[
        \dots
    \]
    i.e. it is reducible, then the irreducible components must be integer spin representations of \(\mathfrak{sl}(2; \mathbb{C})\), which have a \(J_3^{(\alpha)}\)-eigenvalue 0 subspace, which is one dimensional, so that the \(J_3^{(\alpha)}\)-eigenvalue 0 subspace of \(V\) will have dimension \(N\). However from the definition of \(V\) we have that this eigenspace is simply
    \[
        \left\{ \mu J_3^{(\alpha)} \big | \mu \in \mathbb{C} \right\},
    \]
    which is one dimensional, so \(N = 1\) and thus \(V\) is irreducible.
\end{proof}

This leads to a set of practically useful facts about the roots of semi-simple Lie algebras:

\begin{theorem}
    Let \(\mathfrak{g} = \mathfrak{h} \oplus \bigoplus_{\alpha \in \mathcal{R}} \mathfrak{g}_\alpha \) be the root decomposition of a complex semi-simple Lie algebra \(\mathfrak{g}\) with respect to a Cartan subalgebra \(\mathfrak{h}\). Then the following properties hold:
    \begin{enumerate}[label=(\roman*)]
        \item a
        \item Every root subspace is one dimensional, i.e., \(\dim(\mathfrak{g}_\alpha) = 1\) for all \(\alpha \in \mathcal{R}\).
        \item This particular pairing between vectors we are claiming that they are in \(\mathbb{Z}\) i.e. they are integers...
        \item On the dual space we can define a linear map called reflection which takes any vector... if we tae anothr root, its reflection is also a root... if we reflects two roots with respect to each other, we get a rotation in the plane spanned by these two roots, and this rotation maps roots to roots...
        \item For any \(\alpha \in \mathcal{R}\), the only multiples of \(\alpha\) which are roots are \(\pm \alpha\).
        \item For roots \(\alpha,\,\beta\neq\alpha\), then any
              \[
                  W = \bigoplus_{k \in \mathbb{Z}} \mathfrak{g}_{\beta + k \alpha},
              \]
              is an irreducible representation of the root subalgebra \(\mathfrak{sl}(2; \mathbb{C})_{\alpha}\) associated to \(\alpha\).
    \end{enumerate}
\end{theorem}
\begin{proof}
    We can prove these properties one by one.
    \begin{enumerate}[label=(\roman*)]
        \item If \(\mathcal{R} \) does not span \(\mathfrak{h}^{\vee}\) then there exists a non-zero element \(H \in \mathfrak{h}\) such that \(\alpha(H) = 0\) for all \(\alpha \in \mathcal{R}\). This means that \(H\) commutes with all the generators of the algebra, thus it is in the center of the algebra, which contradicts the assumption of semi-simplicity...
        \item This follows from Prop. 3.4.8: since \(\{\mu J_3^{(\alpha)} | \mu \in \mathbb{C} \} \bigoplus\) ...
        \item \(J_3^{(\alpha)}\) every eigenvalue has to be integer or half integer, thus this fraction must be an integer... the pairing is not symmetric and it is only linear in the second argument.
        \item For any pair of roots their angular bracket is integer and from the previous proof we saw that \(E_{\beta} \in \mathfrak{g}_{\beta}\) has \(J_3^{(\alpha)}\) eigenvalue is \(\tfrac{n}{2}\), then there must be.... we can think of it as ....
              \[
                  0 \neq \ket{j,\,-\tfrac{n}{2}} = \begin{dcases}
                      [\dots ] \quad n>0, \\
                      [\dots ] \quad n<0.
                  \end{dcases}
              \]
              Thus in the end we get
              \[
                  \implies \dots
              \]
        \item Assume that \(\alpha\) and \(\beta = c \alpha\), \(c \in \mathbb{C}\), are both roots. By (iii), \(\langle \alpha, \beta \rangle = 2c\) and\( \langle \beta, \alpha \rangle = 2/c \) must be integer, so we have the possibilities \(c = \pm 1, \, \pm 2,\, \pm \tfrac{1}{2}\)...

              By part (ii) of the theorem, we know...

              So \(2 \alpha\) is not a root. Analogously, we can argue for \(\mathfrak{g}_{-\alpha} = \{\mu J_{-}^{(\alpha)} | \mu \in \mathbb{C} \}\) to be the lowest weight subspace, and so \(\mathfrak{g}_{m \alpha} = \{0\}\) for \(m < 1\), i.e., \(-2 \alpha\) is not a root.
        \item ...
    \end{enumerate}
\end{proof}

The next corollary will directly tell us how to construct the entire algebra from the root system, and thus it is the key to the classification of complex semi-simple Lie algebras.

\begin{corollary}
    We only need to know the finite set of root vectors and the pairings among them to reconstruct the entire algebra.

    As a vectors space, the algebra is a direct sum of one dimensional root spaces, each associated to a root, plus the Cartan subalgebra. The Lie bracket is determined by the root structure and the pairing among the root vectors.

    \begin{itemize}
        \item For each root we find basis elements and generators for the cartan subspace which carries some symetric inner product, which along with the pairing directly defines the killing form. The norm is already defined

        \item For every \(X = J_3^{(\alpha)}\), \(E_{\alpha}\) we can define the adjoint action of every generator of the algebra on it (define the crresponding endomorphism \(\mathrm{ad}_X \in \mathrm{End}(\mathfrak{g})\)), which is determined by the root structure and the pairing, thus we can find the commutation relations among all the generators of the algebra:
              \begin{itemize}
                  \item \(\dots \)
                  \item \(\dots \)
                  \item \(\dots \)
              \end{itemize}
    \end{itemize}
\end{corollary}
the last is the only non trivial, but after showing that we have a irreducible representation of \(\mathfrak{sl}(2; \mathbb{C})\) we can use the known representation theory of \(\mathfrak{sl}(2; \mathbb{C})\) so that we can use \(E_{\alpha} \) as a raising operator to find the commutation relations with all the other generators of the algebra, which are determined by the root structure and the pairing...
\begin{proof}
    Almost everything follows straightforwardly from Theorem 3.4.9 and its proof. The only non-trivial points are the numerical factors in the last equation. To see these...
\end{proof}

\subsubsection{Summary}

Let us summarize what we have seen so far.
\begin{itemize}
    \item By viewing a complex semi-simple \(\mathfrak{g}\) as the representation space of the adjoint
          representation (ad, \(\mathfrak{g}\)), we are naturally led to consider a vector space basis for \(\mathfrak{g}\) consisting of simultaneous eigenvectors of the maximally commuting subalgebra (the Cartan subalgebra) \(\mathfrak{h}\) of \(\mathfrak{g}\).
    \item The standard diagonalization theorem directly implies the root decomposition Theorem 3.4.3, where the eigenspaces \(\mathfrak{g}_{\alpha}\) are labelled by the eigenvalues \(\alpha \in\) \(\mathcal{R} \subset \mathfrak{h}^{\vee}\)  called roots.
    \item Every root \(\alpha \in \mathcal{R}\) is associated with a subalgebra \(\mathfrak{sl}(2; \mathbb{C})_{\alpha} \subset \mathfrak{g}\); through the adjoint action, \(\mathfrak{g}\) becomes a (generally reducible) representation of \(\mathfrak{sl}(2; \mathbb{C})_{\alpha}\). Using the known representation theory of \(\mathfrak{sl}(2; \mathbb{C})\), we could deduce a series of properties of \(\mathfrak{g}\) and the roots, which culminated in Theorem 3.4.9.
    \item These results allow us to essentially completely reconstruct the algebra \(\mathfrak{g}\) from its roots. In turn, the stringent consistency conditions of roots make them much easier to quantify and categorize. This will become particularly clear in the following section.
\end{itemize}